% =====================================================================
%  CT Dose AI — the machine learning, in four pages
%  pdflatex CT_Dose_AI_ML_Methods.tex   (twice)
% =====================================================================
\documentclass[10pt,a4paper,twocolumn]{article}

\usepackage[T1]{fontenc}
\usepackage[utf8]{inputenc}
\usepackage{helvet}\renewcommand{\familydefault}{\sfdefault}
\usepackage{microtype,amsmath,amssymb}
\usepackage{flushend}   % balance the two columns on the last page
\usepackage[a4paper,top=16mm,bottom=16mm,left=14mm,right=14mm,columnsep=6mm]{geometry}
\usepackage{titlesec,enumitem,booktabs,tabularx,array,colortbl,xcolor,graphicx}
\usepackage[font=small,labelfont=bf,labelsep=period]{caption}
\usepackage[hidelinks]{hyperref}

\definecolor{ctblue}{HTML}{2A78D6}
\definecolor{ctdark}{HTML}{1A4F94}
\definecolor{ctaqua}{HTML}{1BAF7A}
\definecolor{ctink2}{HTML}{52514E}
\definecolor{ctsurf}{HTML}{F4F4F2}

\titleformat{\section}{\normalsize\bfseries\color{ctdark}}{\thesection\;}{0pt}{}
\titlespacing*{\section}{0pt}{10pt}{3pt}
\titleformat{\subsection}{\small\bfseries\color{ctink2}}{\thesubsection\;}{0pt}{}
\titlespacing*{\subsection}{0pt}{7pt}{2pt}
\setlist[itemize]{leftmargin=11pt,itemsep=1pt,topsep=2pt}
\setlength{\parskip}{3pt}\setlength{\parindent}{0pt}
\renewcommand{\arraystretch}{1.18}
\newcommand{\ctdi}{\ensuremath{\mathrm{CTDI}_{\mathrm{vol}}}}
\newcolumntype{L}{>{\raggedright\arraybackslash}X}
\newcommand{\hdr}{\rowcolor{ctsurf}}

\begin{document}

\twocolumn[{%
\begin{center}
{\LARGE\bfseries Machine learning in CT Dose AI}\\[3pt]
{\large\color{ctink2} Model selection, uncertainty and sequential decision making\\
for radiation dose in coronary CT angiography}\\[6pt]
{\small Thodoris Fragos \quad\textbullet\quad November 2026 \quad\textbullet\quad
\href{https://github.com/Thodoris1657/CT-Dose-AI-Algorithm}{github.com/Thodoris1657/CT-Dose-AI-Algorithm}}
\end{center}
\vspace{4pt}
\hrule
\vspace{8pt}
}]

\section{What the project has to decide}

The dataset is 89 coronary CT angiography examinations with four inputs — BMI, tube voltage, tube
current--time product and a scan-size field — and two dose outputs, \ctdi{} and DLP. The clinical request was
twofold: predict the dose, and say what to change to reduce it. Those are different problems and I treat them
separately throughout.

\textbf{Prediction} is supervised regression on 89 rows. \textbf{Reduction} is a constrained decision problem:
minimising dose alone is degenerate, because the minimum is no radiation at all. The constraint is image
quality, and image quality is absent from the data. Most of the engineering effort went into that gap rather
than into the regression.

\section{The supervised core}

\subsection{Parametrisation}

Dose quantities are positive and their drivers multiply, so every model works on logarithms. CT physics gives
the functional form outright:
\[
\ctdi = k\cdot \mathrm{mAs}\cdot(\mathrm{kV}/120)^{n},\qquad
\mathrm{DLP} = \ctdi \times L .
\]
In log space this is linear, which allows two deliberate restrictions. First, the mAs exponent is \emph{fixed}
at one rather than fitted: dose is linear in tube charge by definition, so fitting that coefficient would spend
a degree of freedom on a known constant and would make extrapolation to unseen mAs unsafe. Second, the fit uses
a Huber loss, because the data contains a minority cluster of examinations with a different dose yield
(\S\ref{sec:quality}) and I did not want seven rows pulling the scanner constant.

\subsection{The candidate set}

I did not assume the physics model was best; it had to win a competition. Nine candidates, all fitted on the
same log targets:

\begin{itemize}
\item the physics law, and physics $+$ BMI, physics $+$ length;
\item \textbf{hybrid residual models} — physics prior plus a learned correction from a Gaussian process
(RBF kernel, white-noise term), a small multilayer perceptron, or histogram gradient boosting;
\item \textbf{purely statistical} — log--log ridge with second-order interactions;
\item \textbf{black box} — histogram gradient boosting with monotonicity constraints on mAs and kV, and a bagged
ensemble of deep multilayer perceptrons (the ``deep learning'' entry).
\end{itemize}

The hybrids are the interesting ones: if the physics law is complete, the residual learner should collapse
towards zero and the hybrid should reduce to the physics law, costing only variance. That is exactly what
happens, and the cost is measurable.

\subsection{Validation protocol}

Repeated grouped $K$-fold cross-validation, $K=5$ with 10 repeats, so each model is fitted 50 times. The
grouping matters: the dataset contains exact duplicate rows, and a hash of the feature--target tuple forces
duplicates into the same fold. Without that, a duplicated patient would appear in training and test
simultaneously and flatter every model.

Selection uses the \textbf{1-standard-error rule} — the simplest model whose mean error is within one standard
error of the best. With 87 training rows the differences between the top models are within noise, and the rule
prevents a complex model winning on a fluctuation.

\subsection{Results}

\begin{table}[h]\footnotesize
\begin{tabularx}{\columnwidth}{@{}L r r@{}}
\toprule
\hdr \textbf{Model} & \textbf{\ctdi} & \textbf{DLP}\\
\midrule
\rowcolor{ctsurf!60}\textbf{Physics law} (selected) & \textbf{0.77\%} & \textbf{9.8\%}\\
Physics $+$ GP residual & 0.97\% & 9.9\%\\
Physics $+$ NN residual & 1.15\% & 9.8\%\\
Log--log ridge (poly2) & 1.08\% & 10.1\%\\
Deep MLP ensemble & 1.31\% & 13.1\%\\
Monotone gradient boosting & 4.53\% & 11.7\%\\
\bottomrule
\end{tabularx}
\caption{Cross-validated mean absolute percentage error, 89 examinations (87 after screening).}
\end{table}

The honest headline is that the flexible models lose, and they lose for a reason that is worth stating: the
functional form is known exactly and the sample is small, so a network spends its capacity re-learning a
straight line from 87 points and does it worse.

\subsection{A prospective test}

Cross-validation on 89 rows can still flatter a model. The dataset arrived in two batches — 50 examinations,
then 39 more — which allows a genuine held-out test. I refitted on the original 50 and predicted the 37 clean
new examinations the model had never seen:

\begin{table}[h]\footnotesize
\begin{tabularx}{\columnwidth}{@{}L r r@{}}
\toprule
\hdr \textbf{Trained on 50, tested on 37} & \textbf{\ctdi} & \textbf{DLP}\\
\midrule
\rowcolor{ctsurf!60}\textbf{Physics law} & \textbf{0.58\%} & 12.7\%\\
Log--log ridge & 1.01\% & 11.6\%\\
Deep MLP ensemble & 1.89\% & 14.2\%\\
Monotone gradient boosting & 4.27\% & 12.5\%\\
\bottomrule
\end{tabularx}
\caption{Median absolute error of the physics law on unseen examinations was 0.15\%.}
\end{table}

The fitted constants moved by almost nothing when the sample grew by 78\%: $k$ from 0.06050 to 0.06048 and $n$
from 2.533 to 2.531. A model whose parameters are stable under a large increase in data is estimating a
property of the scanner rather than fitting the sample, and that is the argument I would defend.

\section{Uncertainty}

A point estimate without an interval invites misuse, so predictions carry \textbf{conformal prediction
intervals} built on out-of-fold log residuals $r_i=\log y_i-\log\hat y_i$. The bounds are empirical quantiles
with the finite-sample correction,
\[
q_{\mathrm{lo}}=Q_{\lfloor (n+1)\alpha/2\rfloor/n}(r),\quad
q_{\mathrm{hi}}=Q_{\lceil (n+1)(1-\alpha/2)\rceil/n}(r),
\]
giving the multiplicative interval $[\hat y e^{q_{\mathrm{lo}}},\,\hat y e^{q_{\mathrm{hi}}}]$. The attraction is
that coverage holds under exchangeability alone — no normality, no homoscedasticity. Measured out-of-fold
coverage is 92\% against a 90\% target for both quantities, and the \ctdi{} interval is visibly asymmetric
($\times[0.93,1.00]$), which is the deviant-yield cluster showing up in the lower tail exactly as it should.

Separately, every prediction is checked against an \textbf{applicability domain}: inputs outside the training
ranges are flagged, and any tube voltage absent from the data (80, 100 kV here) raises an explicit
extrapolation warning in the interface. The model is allowed to extrapolate — the physics supports it — but it
is not allowed to do so silently.

\section{Predicting what the scanner will do}

The dose model needs the mAs, which the automatic tube-current modulation only chooses once the patient is
positioned. A second stage predicts that choice from BMI alone, so the dose can be estimated before the scan.
There is no formula here, and the ranking inverts:

\begin{table}[h]\footnotesize
\begin{tabularx}{\columnwidth}{@{}L r r@{}}
\toprule
\hdr \textbf{From BMI alone} & \textbf{mAs} & \textbf{length}\\
\midrule
Median, ignoring BMI & 30.6\% & 10.5\%\\
Log-linear in BMI & 23.6\% & 5.8\%\\
Log-quadratic & 21.2\% & \textbf{5.0\%}\\
Gaussian process & 21.3\% & 5.3\%\\
Random forest & 23.6\% & 5.5\%\\
\rowcolor{ctsurf!60}\textbf{Neural network} & \textbf{20.8\%} & 5.3\%\\
\bottomrule
\end{tabularx}
\caption{The one task where flexible models win. On 50 examinations a random forest was selected; on 89 the
network takes over, so the ranking is still moving.}
\end{table}

Chaining two models means chaining two error distributions. Rather than adding variances under a normality
assumption, I convolve the two empirical residual distributions by Monte Carlo, which is one line of code and
assumes nothing. The resulting pre-scan interval is roughly $\pm$35\% against $\pm$7\% once the mAs is known.
That width is not a defect; it is the quantitative statement of how unpredictable the modulation is.

\section{Data quality as an algorithm}
\label{sec:quality}

Dose values are typed by hand off a console, and in a dataset of this size two mistyped numbers do real damage.
Rather than eyeballing the file I built a screen that uses physics as the test statistic:
$\ctdi/\mathrm{mAs}$ must equal the scanner constant for that voltage, and $\mathrm{DLP}/\ctdi$ must be a
plausible irradiated length. Robust location and scale (median, MAD) set the thresholds. Rows are labelled
\texttt{ok}, \texttt{review} (a repeatable deviation — a second scan mode) or \texttt{error} (physically
inconsistent in isolation); only \texttt{error} rows are withheld from training, and nothing is deleted from
the file.

The screen found two rows: one whose DLP implies 66 cm of coverage against 23 cm everywhere else, and one whose
dose per mAs is 27\% off the constant. Excluding them moves \ctdi{} error from 1.18\% to 0.77\% and DLP $R^2$
from 0.744 to 0.903. It also identified seven examinations sharing a $-7\%$ dose yield with almost no spread,
which is a second scan mode rather than noise — a finding for the department, not a cleaning step.

A related negative result is worth recording. With 87 clean examinations, the irradiated length implied by
$\mathrm{DLP}/\ctdi$ is \emph{not} recoverable from the four inputs: a random forest achieves a
cross-validated $R^2$ of 0.05. The residual DLP error is therefore irreducible with the current columns, which
turns a suspicion into a measurement and gives the data provider a concrete request.

\section{The decision layer}

Minimising dose subject to constant image quality has a closed form. For a fixed patient, quantum noise scales
as $\sigma\propto D^{-1/2}$ and iodine contrast as $C\propto \mathrm{kV}^{-\gamma}$ with $\gamma\approx1.35$,
so holding contrast-to-noise at a fraction $\tau$ of the current protocol requires
\[
D_{\mathrm{new}}\;\ge\;D_{\mathrm{old}}\,\tau^{2}
\left(\frac{\mathrm{kV_{new}}}{\mathrm{kV_{old}}}\right)^{\!2\gamma}.
\]
The engine evaluates each available voltage, inverts the \ctdi{} model for the compensating mAs, and ranks the
results by DLP. Three constraints are enforced: a minimum voltage by BMI, the tube's maximum output, and the
physical scan length. A fourth comparison is empirical — quantile regression of local dose against BMI gives
the median and 25th-percentile dose actually delivered to similar patients, which is an achievable target that
requires no extrapolation at all.

\section{Acquiring the missing variable}

The constraint above uses a literature exponent. Two active-learning components replace it with measurement.

\subsection{Bayesian optimisation on a phantom}

Objective: minimise \ctdi{} subject to measured CNR $\ge$ target. Dose is known analytically; CNR costs a
phantom scan. The surrogate on $\log\mathrm{CNR}$ is a Bayesian linear model with physics priors
($\gamma\sim\mathcal{N}(1.35,0.5^2)$, noise--dose exponent $\delta\sim\mathcal{N}(0.5,0.15^2)$) plus a Gaussian
process residual for what the formula misses, such as photon starvation at low voltage. The priors are what
make the method usable after two or three scans rather than thirty. Acquisition is constrained expected
improvement,
\[
a(x)=\max\!\left(0,\,d^{*}-d(x)\right)\cdot\Pr\!\left[\mathrm{CNR}(x)\ge\text{target}\right].
\]
In simulation against a synthetic phantom with a known optimum, twelve campaigns converged in a median of eight
scans to within about 8\% of the true constrained optimum, meeting the CNR target in 10 of 10 runs.

\subsection{A safe contextual bandit in routine}

Context is BMI; actions are voltage $\times$ dose-level combinations; the reward is a radiologist's 1--4
diagnostic score, binarised at 3. The model is Bayesian logistic regression fitted by Laplace approximation,
with physics-informed priors (more dose helps, lower voltage helps at equal dose, larger patients are harder).
One shared model means every examination informs every action, which matters enormously at these sample sizes.
Action selection is Thompson sampling: draw one posterior sample, take the lowest-dose action whose sampled
acceptability clears the target. Exploration is therefore proportional to remaining uncertainty and decays on
its own.

Three safety constraints make it deployable rather than merely interesting: the BMI voltage floor; a
pessimistic (10th-percentile) acceptability above 0.85; and a step limit forbidding proposals more than 15\%
below the lowest dose already \emph{proven} acceptable at a similar BMI — the SafeOpt idea of expanding only
into a region the posterior already considers safe. A human approves every proposal. Simulated against a reader
whose acceptance follows true iodine CNR, 1\,500 patients give a 45--50\% dose reduction with a 3--5\%
non-diagnostic rate, consistent with the 0.97 acceptability target that the department sets.

\section*{\small References}
{\footnotesize
[1] Stocker TJ et al. \textit{PROTECTION VI}. Eur Heart J, 2018. ---
[2] Vovk, Gammerman, Shafer. \textit{Algorithmic Learning in a Random World}. Springer. ---
[3] Gelbart, Snoek, Adams. \textit{Bayesian Optimization with Unknown Constraints}. UAI, 2014. ---
[4] Sui, Gotovos, Burdick, Krause. \textit{Safe Exploration for Optimization with Gaussian Processes}
(SafeOpt). ICML, 2015. ---
[5] Chapelle, Li. \textit{An Empirical Evaluation of Thompson Sampling}. NIPS, 2011. ---
[6] Trattner S et al. \textit{Cardiac-specific conversion factors}. JACC Cardiovasc Imaging, 2018. ---
[7] Hastie, Tibshirani, Friedman. \textit{The Elements of Statistical Learning}, §7.10 (the one-standard-error
rule).
\par}

\vspace{4pt}
{\scriptsize\color{ctink2}\hrule\vspace{3pt}
Decision-support software for research use. Not a medical device, not clinically validated. Any protocol change
requires approval by a medical physicist and the responsible radiologist.\par}

\end{document}
